Entanglement between multiple cavities is crucial for quantum computation. A practical way to generate such entanglement is through coupler-mediated cavity-cavity coupling. Flux-tunable devices such as the Superconducting Nonlinear Asymmetric Inductive eLement (SNAIL) \cite{SNAIL}, the Flux-Tunable Transmon (FTT) \cite{squid}, and the Linear INductive Coupler (LINC) \cite{linc} enable beam-splitter interactions through three-wave mixing, providing a practical route to high-fidelity bosonic gates. Activating these interactions often requires biasing the nonlinear element away from its sweet spot, see example discussed in Sec.\ \ref{sec:parametric_two_mode_control}. In such flux bias point, flux noise causes fluctuations of the nonlinear-element frequency. In the dispersive regime, these fluctuations produce Lamb-shift fluctuations, leading to cavity-frequency fluctuations and thus cavity dephasing. This inherited dephasing is especially harmful because it introduces an error channel that directly competes with the favorable photon-loss-dominated bias. As the coupling strength between the cavity and the nonlinear element is increased to speed up gates, this inherited dephasing can become a leading decoherence mechanism, degrading the bias structure that bosonic QEC and dual-rail encodings rely on for their advantage.
% motivate 3wave ratherthan 4wave
A commonly used strategy is to idle the nonlinear element at a flux sweet spot and tune it away from the sweet spot only when operations are needed. In three-dimensional architectures, fast flux tuning can be implemented using magnetic hoses \cite{magnetichose1,magnetichose2} or on-chip flux lines \cite{li2025,valadares2026}. However, these schemes increase hardware complexity and can slow the operation of interest. Moreover, tuning flux might risk unwanted transition. 

This thesis introduces an alternative method to suppress inherited cavity dephasing without changing the dc flux bias. An weak off-resonant drive applied to the nonlinear mode induces a photon-number-dependent AC-Stark shift of the cavity transition frequency. With suitable drive parameters, the fluctuation of this shift due to noise in nonlinear element, cancels the fluctuation from AC stark shift, thereby creating a driven sweet spot. This mechanism differs from the dynamical sweet spots studied for qubits in \cite{PhysRevApplied.15.034065floquetqubit,PhysRevApplied.12.054015dss2}, where a resonant or near-resonant drive on the qubit itself engineers the sweet spot. Here, by contrast, weak off-resonant drive drive acts on the nonlinear element rather than on the protected cavity modes which encode quantum information, and it substantially suppressing inherited dephasing and introduces only a small correction to the cavity decay rate. For a representative parameter set, Monte Carlo simulations show a \(\sim 20\times\) improvement in the cavity dephasing time, from \(\sim 100~\mu\text{s}\) to \(\sim 2~\text{ms}\) with negligible effect on decay time. Our scheme also differs from dynamical decoupling \cite{PhysRev.175.453DD}. Our scheme does not require cavity gate sequences and therefore avoids the \(\sim 1~\mu\text{s}\) overhead of typical cavity gates \cite{smcc-t465snap,eickbusch2022ecd}. 



